\documentclass[12pt,a4paper]{article}

% =========================================================
% PACKAGES
% =========================================================
\usepackage[utf8]{inputenc}
\usepackage[T1]{fontenc}
\usepackage{lmodern}
\usepackage{geometry}
\usepackage{setspace}
\usepackage{graphicx}
\usepackage{booktabs}
\usepackage{array}
\usepackage{float}
\usepackage{amsmath}
\usepackage{hyperref}
\usepackage{xcolor}
\usepackage{enumitem}
\usepackage{caption}
\usepackage{subcaption}
\usepackage{longtable}
\usepackage{fancyhdr}
\usepackage{listings}

% =========================================================
% PAGE SETUP
% =========================================================
\geometry{
    top=1in,
    bottom=1in,
    left=1in,
    right=1in
}

\onehalfspacing

\hypersetup{
    colorlinks=true,
    linkcolor=blue,
    citecolor=blue,
    urlcolor=blue
}

\pagestyle{fancy}
\fancyhf{}
\lhead{ResQAI}
\rhead{Team APEX}
\cfoot{\thepage}

% =========================================================
% CODE STYLE
% =========================================================
\lstset{
    basicstyle=\ttfamily\small,
    breaklines=true,
    frame=single,
    columns=fullflexible
}

% =========================================================
% DOCUMENT
% =========================================================
\begin{document}

% =========================================================
% TITLE PAGE
% =========================================================
\begin{titlepage}
    \centering

    \vspace*{0.8cm}

    {\Large \textbf{TEAM APEX}}\\[0.8cm]

    {\LARGE \textbf{RESQAI}}\\[0.5cm]

    {\Large \textbf{AI-Assisted Disaster Response Through\\
    Humanitarian Message Classification}}\\[1.2cm]

    {\large A Machine Learning Based Approach for\\
    Real-Time Crisis Information Analysis}\\[1.5cm]

    \vfill

    {\large \textbf{PROJECT REPORT}}\\[0.8cm]

    Submitted in partial fulfillment of the requirements\\
    for the academic project

    \vspace{1.2cm}

    {\large \textbf{TEAM MEMBERS}}\\[0.6cm]

    \begin{tabular}{ll}
        \textbf{Dhwani Agarwal} & \textbf{1024031024} \\
        \textbf{Sanchita Jain} & \textbf{1024031020} \\
        \textbf{Yuvakshi Sood} & \textbf{1024031025}
    \end{tabular}

    \vspace{1.5cm}

    {\large \textbf{DEPARTMENT OF COMPUTER SCIENCE}}\\[0.4cm]

    {\large \textbf{THAPAR INSTITUTE OF ENGINEERING AND TECHNOLOGY}}\\[0.8cm]

    {\large 2026}

\end{titlepage}

% =========================================================
% ABSTRACT
% =========================================================
\begin{abstract}

Natural disasters and humanitarian emergencies generate large amounts of
information through social media and other online communication channels.
During such events, affected individuals may report injuries, request
assistance, provide information about infrastructure damage, communicate
evacuation requirements, or offer rescue and volunteering support. However,
these useful messages are mixed with irrelevant, political, emotional, and
general informational content. Manually filtering this information is
time-consuming and difficult during rapidly evolving emergencies.

This project presents \textbf{ResQAI}, an AI-assisted disaster information
processing system designed to automatically classify crisis-related
social-media messages into humanitarian categories. The system uses
\textbf{BERTweet}, a transformer-based language model pretrained specifically
for Twitter-style text, and fine-tunes it on the \textbf{HumAID} dataset.

Multiple approaches were investigated, including traditional TF-IDF-based
machine learning models, Logistic Regression, Linear Support Vector Machines,
Random Forest, DistilBERT, and BERTweet. Among the evaluated approaches,
BERTweet produced the strongest overall performance. The final BERTweet model
achieved \textbf{78.67\% accuracy} and a \textbf{78.27\% weighted F1 score}
on the unseen test set.

Error analysis revealed that several humanitarian categories have substantial
semantic overlap. In particular, the
\texttt{other\_relevant\_information} category was frequently confused with
categories such as \texttt{caution\_and\_advice} and
\texttt{not\_humanitarian}. To address the limitations of a single
classification model, a hybrid architecture is proposed in which BERTweet
acts as the primary fast classifier while a large language model such as
Gemini can provide secondary contextual analysis for low-confidence messages.

The proposed system is intended as a decision-support and information
filtering system rather than an autonomous emergency-response authority.

\end{abstract}

\newpage

% =========================================================
% TABLE OF CONTENTS
% =========================================================
\tableofcontents
\newpage

% =========================================================
% 1. INTRODUCTION
% =========================================================
\section{Introduction}

Natural disasters such as earthquakes, floods, hurricanes, wildfires, and
other humanitarian emergencies can create rapidly changing situations in
which timely information is extremely valuable. During such events, people
frequently use social media to report incidents, request assistance, share
warnings, communicate their condition, or provide information about affected
areas.

Although social media can provide valuable real-time situational information,
the sheer volume of messages makes manual processing impractical. A stream of
disaster-related messages may contain emergency requests alongside political
discussions, expressions of sympathy, general news, jokes, opinions, and
unrelated content.

ResQAI addresses this problem by automatically classifying incoming messages
into humanitarian information categories. The purpose of the system is to
filter and organize large amounts of information so that potentially useful
messages can be identified more efficiently.

\subsection{Problem Statement}

Given a social-media message associated with a disaster or humanitarian
event, the system should determine the most relevant humanitarian category
of the message.

The classification categories include:

\begin{itemize}
    \item reports of injured or dead people,
    \item infrastructure and utility damage,
    \item requests for urgent assistance,
    \item displaced people and evacuation information,
    \item missing or found people,
    \item rescue, volunteering, and donation efforts,
    \item caution and safety information,
    \item sympathy and support,
    \item other relevant disaster information, and
    \item non-humanitarian information.
\end{itemize}

\subsection{Objectives}

The primary objectives of ResQAI are:

\begin{enumerate}
    \item To develop an automated classifier for humanitarian disaster
    messages.
    
    \item To investigate both classical machine learning and transformer-based
    approaches.
    
    \item To fine-tune a social-media-specific transformer model for disaster
    message classification.
    
    \item To evaluate model performance using validation and unseen test data.
    
    \item To identify and analyze major sources of classification errors.
    
    \item To design a confidence-based architecture capable of using an LLM
    for ambiguous messages.
\end{enumerate}

\subsection{Motivation}

During a disaster, emergency responders may receive thousands of messages
within a short period of time. Not every message requires immediate action,
and important messages can be buried among large quantities of less useful
content.

An automated classification layer can therefore act as an initial information
filter. Rather than replacing emergency personnel, such a system can help
organize incoming information and highlight potentially relevant categories.

% =========================================================
% 2. BACKGROUND
% =========================================================
\section{Background and Related Technology}

\subsection{Natural Language Processing}

Natural Language Processing (NLP) is a field of artificial intelligence
concerned with enabling computers to process and analyze human language.

Social-media text presents several challenges compared with formal written
language. These include:

\begin{itemize}
    \item abbreviations and informal expressions,
    \item spelling variations,
    \item hashtags and mentions,
    \item emojis and emoticons,
    \item incomplete sentences,
    \item noisy punctuation, and
    \item context-dependent meanings.
\end{itemize}

These characteristics make social-media-specific language models particularly
useful for this application.

\subsection{TF-IDF}

Traditional machine learning experiments in this project used Term
Frequency--Inverse Document Frequency (TF-IDF) to convert text into numerical
vectors.

The TF-IDF value of term $t$ in document $d$ can be expressed as:

\begin{equation}
TFIDF(t,d) = TF(t,d) \times IDF(t)
\end{equation}

where:

\begin{equation}
IDF(t) = \log\left(\frac{N}{df(t)}\right)
\end{equation}

Here, $N$ is the total number of documents and $df(t)$ is the number of
documents containing term $t$.

TF-IDF provides an efficient baseline representation but does not directly
model contextual relationships between words.

\subsection{Transformer Models}

Transformer architectures use self-attention mechanisms to capture contextual
relationships between tokens. Transformer-based language models can therefore
produce contextual representations that are more suitable for understanding
natural language than simple frequency-based representations.

\subsection{BERTweet}

BERTweet is a transformer-based language model pretrained specifically on
English Twitter data. Since ResQAI targets short and informal social-media
messages, BERTweet is particularly suitable for the application.

The pretrained BERTweet model was fine-tuned using the HumAID dataset for
ten-class humanitarian classification.

% =========================================================
% 3. DATASET
% =========================================================
\section{Dataset}

The project uses the \textbf{HumAID} dataset, a humanitarian information
dataset containing disaster-related social-media messages.

\subsection{Dataset Distribution}

The dataset used in the experiments consisted of:

\begin{table}[H]
\centering
\caption{HumAID Dataset Distribution}
\begin{tabular}{lr}
\toprule
\textbf{Dataset Split} & \textbf{Number of Messages} \\
\midrule
Training & 53,531 \\
Validation & 7,793 \\
Testing & 15,160 \\
\midrule
\textbf{Total} & \textbf{76,484} \\
\bottomrule
\end{tabular}
\end{table}

\subsection{Classification Categories}

The dataset contains ten humanitarian information categories.

\begin{longtable}{p{0.48\textwidth}p{0.42\textwidth}}
\caption{HumAID Classification Categories} \\
\toprule
\textbf{Class} & \textbf{Description} \\
\midrule
\endfirsthead

\toprule
\textbf{Class} & \textbf{Description} \\
\midrule
\endhead

\texttt{caution\_and\_advice}
&
Safety warnings and recommendations
\\

\texttt{displaced\_people\_and\_evacuations}
&
Information about displaced people and evacuation
\\

\texttt{infrastructure\_and\_utility\_damage}
&
Reports of damage to buildings, roads, utilities, and other infrastructure
\\

\texttt{injured\_or\_dead\_people}
&
Reports concerning injured or deceased individuals
\\

\texttt{missing\_or\_found\_people}
&
Reports of missing or found individuals
\\

\texttt{not\_humanitarian}
&
Disaster-related content that is not directly humanitarian
\\

\texttt{other\_relevant\_information}
&
Other relevant information concerning the disaster
\\

\texttt{requests\_or\_urgent\_needs}
&
Requests for emergency assistance, resources, or supplies
\\

\texttt{rescue\_volunteering\_or\_donation\_effort}
&
Rescue activities, volunteer efforts, and donation-related information
\\

\texttt{sympathy\_and\_support}
&
Messages expressing sympathy, solidarity, or emotional support
\\

\bottomrule
\end{longtable}

% =========================================================
% 4. METHODOLOGY
% =========================================================
\section{Methodology}

The project followed an iterative machine learning development process.

\subsection{Data Preparation}

The dataset was loaded and divided into training, validation, and test sets.
The class names were converted into numerical identifiers for model training.

Each class was assigned a unique integer:

\begin{equation}
label \rightarrow \{0,1,\ldots,9\}
\end{equation}

The text was then processed using the BERTweet tokenizer.

\subsection{Tokenization}

BERTweet tokenization converts each message into a sequence of token IDs and
an associated attention mask.

The maximum sequence length was set to 128 tokens.

Before selecting this value, token-length statistics were examined. The
average training and validation message lengths were approximately 35 tokens.
Only approximately 0.10\% of training messages and 0.04\% of validation
messages reached the 128-token limit.

Therefore, increasing the maximum sequence length was not expected to provide
a meaningful improvement and would have increased computational requirements.

\subsection{Model Training}

The final BERTweet model was fine-tuned using the following configuration.

\begin{table}[H]
\centering
\caption{Final BERTweet Training Configuration}
\begin{tabular}{ll}
\toprule
\textbf{Parameter} & \textbf{Value} \\
\midrule
Base model & BERTweet-base \\
Maximum sequence length & 128 \\
Learning rate & $1\times10^{-5}$ \\
Training batch size & 8 \\
Evaluation batch size & 8 \\
Maximum epochs & 3 \\
Weight decay & 0.01 \\
Warm-up steps & 500 \\
Learning-rate scheduler & Linear \\
Mixed precision & FP16 \\
Model selection metric & Weighted F1 \\
Early stopping patience & 1 \\
\bottomrule
\end{tabular}
\end{table}

\subsection{Model Selection}

The model was evaluated at the end of every epoch. The checkpoint with the
highest validation weighted F1 score was retained.

This approach is important because lower training loss does not necessarily
mean better generalization. During experimentation, validation performance
stopped improving after approximately the second epoch while training loss
continued to decrease.

The use of early stopping and best-checkpoint selection therefore reduced the
risk of retaining an overfit model.

% =========================================================
% 5. MODEL EXPERIMENTS
% =========================================================
\section{Model Experiments}

Several approaches were evaluated before selecting BERTweet.

\subsection{Traditional Machine Learning Models}

TF-IDF representations were used with classical classifiers including:

\begin{itemize}
    \item Naive Bayes,
    \item Logistic Regression,
    \item class-weighted Logistic Regression,
    \item Linear Support Vector Machine, and
    \item Random Forest.
\end{itemize}

These models provided useful baselines but generally performed below the
transformer-based models.

\subsection{Transformer Models}

Two transformer architectures were investigated:

\begin{itemize}
    \item DistilBERT,
    \item BERTweet.
\end{itemize}

Class weighting was also investigated for DistilBERT and BERTweet-related
experiments. However, class weighting did not improve the final performance
and was therefore not used in the selected BERTweet configuration.

\subsection{Model Comparison}

The approximate experimental results are summarized below.

\begin{table}[H]
\centering
\caption{Comparison of Evaluated Approaches}
\begin{tabular}{lcc}
\toprule
\textbf{Model} & \textbf{Approx. Accuracy} & \textbf{Observation} \\
\midrule
Naive Bayes & $\sim$68--69\% & Baseline \\
Logistic Regression & $\sim$73--75\% & Strong classical baseline \\
Weighted Logistic Regression & $\sim$73\% & No major improvement \\
Linear SVM & $\sim$72--73\% & Competitive classical model \\
Random Forest & $\sim$71--72\% & Lower performance \\
DistilBERT & $\sim$77\% & Strong transformer baseline \\
Weighted DistilBERT & $\sim$76\% & Weighting reduced performance \\
BERTweet & \textbf{$\sim$79\%} & Best overall approach \\
\bottomrule
\end{tabular}
\end{table}

The results indicate that transformer-based models were better suited to
capturing the contextual characteristics of disaster-related social-media
messages.

% =========================================================
% 6. RESULTS
% =========================================================
\section{Results}

\subsection{Validation Performance}

The selected BERTweet configuration showed the following validation behavior:

\begin{table}[H]
\centering
\caption{BERTweet Validation Performance by Epoch}
\begin{tabular}{lcc}
\toprule
\textbf{Epoch} & \textbf{Validation Accuracy} &
\textbf{Weighted F1} \\
\midrule
1 & 76.90\% & 75.64\% \\
2 & \textbf{78.40\%} & \textbf{78.05\%} \\
3 & 78.28\% & 77.78\% \\
\bottomrule
\end{tabular}
\end{table}

The second epoch achieved the highest validation weighted F1 score.

The third epoch showed a small decline in validation F1 despite continued
improvement in training loss. This indicates the beginning of overfitting and
supports the use of best-checkpoint selection.

\subsection{Final Test Performance}

After model selection, the model was evaluated on the unseen test set.

\begin{table}[H]
\centering
\caption{Final BERTweet Test Results}
\begin{tabular}{lc}
\toprule
\textbf{Metric} & \textbf{Result} \\
\midrule
Test Accuracy & \textbf{78.67\%} \\
Weighted F1 Score & \textbf{78.27\%} \\
\bottomrule
\end{tabular}
\end{table}

The relatively close validation and test performance indicates that the
selected model generalizes reasonably well to unseen data.

% =========================================================
% 7. ERROR ANALYSIS
% =========================================================
\section{Error Analysis}

A confusion matrix was used to examine the behavior of the final BERTweet
classifier.

The most significant source of confusion was the
\texttt{other\_relevant\_information} category. Messages belonging to this
class were frequently predicted as other disaster-related categories.

This is understandable because the class itself acts as a broad category for
relevant information that does not clearly belong to another humanitarian
class.

\subsection{Semantic Overlap}

Examples of ambiguous messages include:

\begin{quote}
``Aftershocks pls \#eqnz''
\end{quote}

and:

\begin{quote}
``When will it be safe to drink water, have kids walk around barefoot?''
\end{quote}

These messages contain useful disaster-related information, but their exact
category can depend on context.

For example, an aftershock-related message may be interpreted as:

\begin{itemize}
    \item other relevant information,
    \item caution and advice, or
    \item an urgent situation.
\end{itemize}

Similarly, a question about drinking water may represent a safety warning,
an information request, or general disaster information.

\subsection{Main Sources of Error}

The primary sources of confusion observed during analysis were:

\begin{enumerate}
    \item \texttt{other\_relevant\_information} versus
    \texttt{not\_humanitarian},
    
    \item \texttt{other\_relevant\_information} versus
    \texttt{caution\_and\_advice},
    
    \item \texttt{requests\_or\_urgent\_needs} versus
    \texttt{rescue\_volunteering\_or\_donation\_effort}.
\end{enumerate}

These errors indicate that the performance limitation is not necessarily
caused by insufficient model capacity. The category definitions themselves
have overlapping semantic boundaries.

% =========================================================
% 8. RESQAI SYSTEM ARCHITECTURE
% =========================================================
\section{Proposed ResQAI Architecture}

The final ResQAI concept extends the trained BERTweet classifier with a
confidence-based secondary analysis layer.

\subsection{Primary Classifier}

BERTweet acts as the primary classification model.

For an incoming message, the model produces:

\begin{itemize}
    \item predicted humanitarian category,
    \item prediction confidence.
\end{itemize}

High-confidence predictions can be processed directly by the system.

\subsection{LLM Fallback}

Messages for which BERTweet produces low-confidence predictions can be sent to
a large language model such as Gemini for additional contextual analysis.

The LLM can be instructed to return structured information such as:

\begin{itemize}
    \item humanitarian category,
    \item priority,
    \item location,
    \item number of affected people,
    \item type of assistance required,
    \item recommended action,
    \item explanation.
\end{itemize}

This component is proposed as a secondary reasoning layer rather than as a
replacement for the trained classifier.

\subsection{Hybrid Architecture}

The proposed architecture is:

\begin{verbatim}
                  Incoming Message
                         |
                         v
                 +---------------+
                 |    BERTweet   |
                 |  Classifier   |
                 +-------+-------+
                         |
                 Category + Confidence
                         |
               +---------+---------+
               |                   |
       High Confidence       Low Confidence
               |                   |
               v                   v
        BERTweet Result        Gemini API
                                   |
                                   v
                           Contextual Analysis
                                   |
                                   v
                         Structured Information
                                   |
                                   v
                           Priority / Action
\end{verbatim}

The architecture allows the domain-specific classifier to handle clear
messages efficiently while providing additional reasoning for ambiguous
messages.

% =========================================================
% 9. SYSTEM WORKFLOW
% =========================================================
\section{System Workflow}

The proposed workflow is as follows:

\begin{enumerate}
    \item Receive an incoming disaster-related social-media message.
    
    \item Tokenize the message using the BERTweet tokenizer.
    
    \item Run the trained BERTweet classifier.
    
    \item Obtain the predicted humanitarian category and confidence score.
    
    \item Compare the confidence score with a predefined threshold.
    
    \item If the confidence is sufficiently high, use the BERTweet prediction.
    
    \item If the confidence is low, send the message to the secondary LLM
    analysis layer.
    
    \item Extract structured information such as location, urgency, and
    assistance requirements.
    
    \item Present the resulting information to the user or emergency
    information interface.
\end{enumerate}

% =========================================================
% 10. EXAMPLE SYSTEM OPERATION
% =========================================================
\section{Example System Operation}

Consider the following incoming message:

\begin{quote}
``People are trapped inside the collapsed building and urgently need rescue.''
\end{quote}

The BERTweet classifier may identify the message as:

\begin{verbatim}
Category: injured_or_dead_people
Confidence: High
\end{verbatim}

Since the message is classified with high confidence, the system can use the
BERTweet prediction directly.

Now consider a more ambiguous message:

\begin{quote}
``Not sure if I slept through the aftershocks or if they are starting again.''
\end{quote}

If the classifier produces a low-confidence prediction, the system can invoke
the secondary LLM layer. The LLM can provide additional contextual
interpretation and structured information such as the likely category,
location, urgency, and recommended action.

This demonstrates why the proposed architecture does not simply send every
message directly to an LLM.

% =========================================================
% 11. ADVANTAGES
% =========================================================
\section{Advantages}

The proposed ResQAI architecture provides several advantages.

\subsection{Domain-Specific Classification}

BERTweet is fine-tuned specifically on humanitarian disaster messages rather
than being used as a generic text classifier.

\subsection{Social-Media Suitability}

BERTweet was pretrained on Twitter-style data, making it appropriate for
informal and short messages.

\subsection{Reduced LLM Dependence}

High-confidence messages can be processed directly by the trained classifier,
reducing unnecessary external LLM calls.

\subsection{Confidence-Based Reasoning}

The system can reserve deeper LLM analysis for ambiguous messages.

\subsection{Scalability}

Automated classification can process large volumes of messages more
efficiently than manual inspection.

\subsection{Decision Support}

The system is designed to organize information and support human decision
making rather than independently making emergency decisions.

% =========================================================
% 12. LIMITATIONS
% =========================================================
\section{Limitations}

The current implementation has several limitations.

\begin{enumerate}
    \item The final classifier achieves approximately 79\% accuracy rather
    than perfect classification.
    
    \item Several humanitarian categories have overlapping semantic meanings.
    
    \item Social-media language is highly informal and context-dependent.
    
    \item The model is primarily trained on English-language data.
    
    \item A confidence threshold cannot guarantee that all ambiguous messages
    will be correctly interpreted.
    
    \item The proposed LLM fallback architecture requires external API access.
    
    \item Real-world emergency deployment would require extensive reliability,
    security, privacy, and human-oversight testing.
\end{enumerate}

Therefore, ResQAI should be treated as an information filtering and
decision-support system rather than an autonomous emergency-response system.

% =========================================================
% 13. FUTURE SCOPE
% =========================================================
\section{Future Scope}

Several improvements can be explored in future versions of ResQAI.

\subsection{Multilingual Support}

Disaster communication may occur in multiple languages. Multilingual
transformer models could be introduced to support regional and international
disaster communication.

\subsection{Location Extraction}

A dedicated entity extraction component could identify locations mentioned in
messages and map them to geographical coordinates.

\subsection{Real-Time Message Streams}

The classifier could be integrated with real-time social-media or emergency
message streams to support continuous monitoring.

\subsection{Priority Prediction}

A separate model could predict the urgency of a message independently from
the humanitarian category.

\subsection{Multimodal Analysis}

Images and videos shared during disasters could be analyzed together with
textual information.

\subsection{Human-in-the-Loop Verification}

Messages identified as critical could be presented to trained emergency
operators for verification before any operational action is taken.

\subsection{Improved Classification Boundaries}

Further work could investigate hierarchical classification or multilabel
classification because some disaster messages naturally contain information
belonging to multiple categories.

% =========================================================
% 14. CONCLUSION
% =========================================================
\section{Conclusion}

This project presented \textbf{ResQAI}, an AI-assisted system for processing
and classifying humanitarian disaster messages.

Multiple machine learning and transformer-based approaches were evaluated.
Traditional TF-IDF-based approaches provided useful baselines, while
transformer-based approaches produced stronger performance. BERTweet was
selected as the primary model because of its suitability for social-media
language and its superior performance in the experiments.

The final BERTweet model achieved:

\begin{center}
\Large
\textbf{78.67\% Test Accuracy}\\[0.2cm]
\textbf{78.27\% Weighted F1 Score}
\end{center}

on the unseen test set containing 15,160 messages.

Error analysis showed that the largest challenge was semantic overlap between
humanitarian categories, particularly the broad
\texttt{other\_relevant\_information} category and related classes.

The project therefore proposes a hybrid architecture in which BERTweet acts
as the primary domain-specific classifier and a large language model can be
used as a secondary reasoning layer for low-confidence messages. This
architecture combines the efficiency and measurable performance of a
specialized classifier with the contextual reasoning capabilities of an LLM.

ResQAI is ultimately intended to help organize the large volume of
information generated during disasters and surface potentially important
information for further human assessment.

% =========================================================
% REFERENCES
% =========================================================
\newpage

\begin{thebibliography}{9}

\bibitem{bertweet}
D. Q. Nguyen, T. Vu, and A. T. Nguyen,
``BERTweet: A pre-trained language model for English Tweets,''
in \textit{Proceedings of the 2020 Conference on Empirical Methods in Natural
Language Processing: System Demonstrations}, 2020.

\bibitem{bert}
J. Devlin, M.-W. Chang, K. Lee, and K. Toutanova,
``BERT: Pre-training of Deep Bidirectional Transformers for Language
Understanding,''
in \textit{Proceedings of NAACL-HLT}, 2019.

\bibitem{transformer}
A. Vaswani et al.,
``Attention Is All You Need,''
in \textit{Advances in Neural Information Processing Systems}, 2017.

\bibitem{tfidf}
G. Salton and C. Buckley,
``Term-weighting approaches in automatic text retrieval,''
\textit{Information Processing \& Management},
vol. 24, no. 5, pp. 513--523, 1988.

\bibitem{humaid}
QCRI,
``HumAID: Humanitarian Information Dataset,''
Humanitarian information classification dataset.

\end{thebibliography}

% =========================================================
% END
% =========================================================

\end{document}